\begin{tcolorbox}[colback=red!3!white,colframe=red!50!black,boxrule=0.8pt,title={Correction notes},fonttitle=\bfseries]
\textbf{Corrections from the calculation review (30~Sept.~2026).} The following places are corrected or annotated in the text; each carries a dated note there with the previous value:
\begin{itemize}
\item $\Lambda_{\text{T0}} = E_P/\xi$: $7.5 \times 10^{15} \to 9.2 \times 10^{22}$~GeV.
\item Hierarchy: $1/\sqrt\xi = 86.6$ calculated correctly, measured $E_P/v = 5 \times 10^{16}$ -- note.
\item Hydrogen correction: $\sim 10^{-32} \to \sim 2 \times 10^{-30}$~eV; atom interferometer: $10^{-18} \to 2 \times 10^{-23}$--$2 \times 10^{-22}$~rad ($v = 1$--$10$~m/s); likewise in the summaries.
\item Addendum 1~Oct.~2026: spinor connection $\Gamma_\mu^{(T)}$: $-\partial_\mu\deltaE/\deltaE^2 \to -\partial_\mu\deltaE/\deltaE$; T0 potential $\xi\hbar^2\,\deltaE/E_{\text{Pl}}$ (dimensionally wrong) $\to \xi\,\deltaE^2/E_{\text{Pl}}$.
\end{itemize}
\end{tcolorbox}

\section*{Abstract}
		This presentation outlines how central aspects of quantum field theory, quantum mechanics, and quantum computer technology can be formulated within the T0-Framework. Based on the time-mass duality $T_{\text{field}} \cdot \Efield = 1$ and the universal parameter $\xipar = \frac{4}{3} \times 10^{-4}$, extensions of the Schrödinger and Dirac equations, corrections to the Bell inequalities, and concepts for deterministic quantum computers are proposed. The document presents a deterministic, local-realistic interpretation that offers a route to solving the measurement problem of quantum mechanics, and it outlines possible applications in quantum technology. Many of these approaches are not yet derived in detail in this early document and are therefore marked as speculative.

	\medskip\noindent\textbf{Note on versions.} The German version of this document is more extensive than this English one. Earlier revisions shortened or changed parts of the English version. Where the two differ, the more extensive German version applies.


	\medskip\noindent\textbf{Status markers.} Statements marked \textbf{[S]} are \emph{speculative}: a hint or conjecture that has not yet been derived or numerically checked in the corpus.

	\section{Introduction: T0 Approach to QFT and QM}
	The T0-Theory proposes extensions of quantum field theory and of the basic equations of quantum mechanics, and it opens new approaches for quantum computer technologies.
	\begin{tcolorbox}[colback=blue!5!white,colframe=blue!75!black,title={T0 Basic Principles for QFT and QM}]
		\textbf{Fundamental T0 Relations:}
		\begin{align}
			T_{\text{field}}(x,t) \cdot \Efield(x,t) &= 1 \quad \text{(Time-Energy Duality)} \\
			\square \deltaE + \xipar \cdot \mathcal{F}[\deltaE] &= 0 \quad \text{(Universal Field Equation)} \\
			\mathcal{L} &= \frac{\xipar}{\EPlanck^2} (\partial \deltaE)^2 \quad \text{(T0 Lagrangian Density)}
		\end{align}
	\end{tcolorbox}
	\section{T0 Field Quantization}
	\subsection{Canonical Quantization with Dynamic Time}
	The central step of T0-QFT lies in the treatment of time as a dynamic field:
	\begin{tcolorbox}[colback=green!5!white,colframe=green!75!black,title={T0 Canonical Quantization}]
		\textbf{Modified Canonical Commutation Relations:}
		\begin{align}
			[\hat{\phi}(x), \hat{\pi}(y)] &= i\hbar \delta^3(x-y) \cdot T_{\text{field}}(x,t) \\
			[\hat{\Efield}(x), \hat{\Pi}_E(y)] &= i\hbar \delta^3(x-y) \cdot \frac{\xipar}{\EPlanck^2}
		\end{align}
	\end{tcolorbox}
	The field operators take an extended form:
	\begin{equation}
		\hat{\phi}(x,t) = \int \frac{d^3k}{(2\pi)^3} \frac{1}{\sqrt{2\omega_k \cdot T_{\text{field}}(t)}} \left[\hat{a}_k e^{-ik \cdot x} + \hat{b}^\dagger_k e^{ik \cdot x}\right]
	\end{equation}
	\subsection{T0-Modified Dispersion Relation}
	The energy-momentum relation is modified by the time field:
	\begin{equation}
		\boxed{\omega_k = \sqrt{k^2 + m^2} \cdot \left(1 + \xipar \cdot \frac{\langle\deltaE\rangle}{\EPlanck}\right)}
	\end{equation}
	\section{T0 Renormalization: Natural Cutoff}
	\begin{tcolorbox}[colback=red!5!white,colframe=red!75!black,title={T0 Renormalization}]
		\textbf{Natural UV-Cutoff:}
		\begin{equation}
			\Lambda_{\text{T0}} = \frac{\EPlanck}{\xipar} \approx 9.2 \times 10^{22} \text{ GeV}
		\end{equation}

		\textit{(Corrected on 30~Sept.~2026: previously $7.5 \times 10^{15}$~GeV -- $1.221 \times 10^{19}\,\text{GeV} \times 7500 = 9.2 \times 10^{22}$~GeV.)}
		At this scale the loop integrals are naturally cut off; a detailed proof of convergence is not given here.
	\end{tcolorbox}
	The beta functions are modified by T0 corrections:
	\begin{equation}
		\beta_g^{\text{T0}} = \beta_g^{\text{SM}} + \xipar \cdot \frac{g^3}{(4\pi)^2} \cdot f_{\text{T0}}(g)
	\end{equation}
	\section{T0 Quantum Mechanics: Basic Equations in the T0 Framework}
	\subsection{T0-Modified Schrödinger Equation}
	The Schrödinger equation receives an extension through the dynamic time field:
	\begin{tcolorbox}[colback=cyan!5!white,colframe=cyan!75!black,title={T0 Schrödinger Equation}]
		\textbf{Time Field-Dependent Schrödinger Equation:}
		\begin{equation}
			i\hbar \cdot T_{\text{field}}(x,t) \frac{\partial\psi}{\partial t} = \hat{H}_0 \psi + \hat{V}_{\text{T0}}(x,t) \psi
		\end{equation}
		where:
		\begin{align}
			\hat{H}_0 &= -\frac{\hbar^2}{2m} \nabla^2 + V_{\text{extern}}(x) \\
			\hat{V}_{\text{T0}}(x,t) &= \xipar \cdot \frac{\deltaE(x,t)^2}{E_{\text{Pl}}}
		\end{align}
		\textit{(Corrected on 1~Oct.~2026: previously $\hat{V}_{\text{T0}} = \xi \hbar^2 \cdot \deltaE/E_{\text{Pl}}$ -- $\xi$ and $\deltaE/E_{\text{Pl}}$ are dimensionless, so the old form has the dimension of $\hbar^2$ and is not an energy like $\hat{H}_0$; a dimensionally correct form is $\xi\,\deltaE^2/E_{\text{Pl}}$. The document does not fix uniquely which normalisation was intended.)}
	\end{tcolorbox}
	\subsubsection{Physical Interpretation}
	The T0 modification leads to three changes:
	\begin{enumerate}
		\item \textbf{Variable Time Evolution:} The quantum evolution proceeds more slowly in regions of high energy density
		\item \textbf{Energy Field Coupling:} The T0 potential couples quantum particles to local field fluctuations
		\item \textbf{Deterministic Corrections:} Subtle, but measurable deviations from standard QM predictions
	\end{enumerate}
	\subsubsection{Hydrogen Atom with T0 Corrections}
	For the hydrogen atom, the result is:
	\begin{align}
		E_n^{\text{T0}} &= E_n^{\text{Bohr}} \left(1 + \xipar \frac{E_n}{\EPlanck}\right) \\
		&= -13.6 \text{ eV} \cdot \frac{1}{n^2} \left(1 + \xipar \frac{13.6 \text{ eV}}{1.22 \times 10^{19} \text{ GeV}}\right)
	\end{align}
	The correction is tiny (relative $1.5 \times 10^{-31}$, absolute $\sim 2 \times 10^{-30}$ eV) \textit{(Corrected on 30~Sept.~2026: previously \enquote{$\sim 10^{-32}$ eV} -- $\xi \cdot 13.6\,\text{eV}/1.22 \times 10^{28}\,\text{eV} = 1.5 \times 10^{-31}$, times $13.6$~eV gives $2 \times 10^{-30}$~eV.)} and lies far beyond present spectroscopic resolution; it is measurable only in principle.
	\subsection{T0-Modified Dirac Equation}
	Relativistic quantum mechanics is extended by the T0 time field:
	\begin{tcolorbox}[colback=magenta!5!white,colframe=magenta!75!black,title={T0 Dirac Equation}]
		\textbf{Time Field-Dependent Dirac Equation:}
		\begin{equation}
			\left[i\gamma^\mu \left(\partial_\mu + \frac{\xipar}{\EPlanck} \Gamma_\mu^{(T)}\right) - m\right]\psi = 0
		\end{equation}
		where the T0 spinor connection is:
		\begin{equation}
			\Gamma_\mu^{(T)} = \frac{1}{\Tfield(x)} \partial_\mu \Tfield(x) = -\frac{\partial_\mu \deltaE}{\deltaE} = -\partial_\mu \ln \deltaE
		\end{equation}
		\textit{(Corrected on 1~Oct.~2026: previously $-\partial_\mu \deltaE/\deltaE^2$ -- with $T \cdot \deltaE = 1$ one has $\partial_\mu T = -\partial_\mu \deltaE/\deltaE^2$; the logarithmic derivative $(1/T)\,\partial_\mu T$ gives $\deltaE \cdot (-\partial_\mu \deltaE/\deltaE^2) = -\partial_\mu \deltaE/\deltaE$.)}
	\end{tcolorbox}
	\subsubsection{Spin and T0 Fields}
	The spin properties are modified by the time field:
	\begin{align}
		\vec{S}^{\text{T0}} &= \vec{S}^{\text{Standard}} \left(1 + \xipar \frac{\langle\deltaE\rangle}{\EPlanck}\right) \\
		g_{\text{factor}}^{\text{T0}} &= 2 + \xipar \frac{m^2}{M_{\text{Pl}}^2}
	\end{align}
	\textbf{[S]} A contribution of this term to the anomalous magnetic moments of the electron and muon is not shown by this; because of the factor $m^2/M_{\text{Pl}}^2$ it is very small for both particles.
	\section{T0 Quantum Computers: Approaches to Information Processing}
	\subsection{Deterministic Quantum Logic}
	The T0 theory suggests a different concept of quantum computers \textbf{[S]}:
	\begin{tcolorbox}[colback=yellow!5!white,colframe=yellow!75!black,title={T0 Quantum Computer Principles}]
		\textbf{Fundamental Differences from Standard QC:}
		\begin{itemize}
			\item \textbf{Deterministic Evolution:} Quantum gates are predictable in principle
			\item \textbf{Energy Field-Based Qubits:} $|0\rangle$, $|1\rangle$ as energy field configurations
			\item \textbf{Time Field Control:} Manipulation through local time field modulation
			\item \textbf{Natural Error Correction:} Self-stabilizing energy fields
		\end{itemize}
	\end{tcolorbox}
	\subsection{T0 Qubit Representation}
	A T0 qubit is realized through energy field configurations:
	\begin{align}
		|0\rangle_{\text{T0}} &\leftrightarrow \deltaE_0(x,t) = E_0 \cdot f_0(x,t) \\
		|1\rangle_{\text{T0}} &\leftrightarrow \deltaE_1(x,t) = E_1 \cdot f_1(x,t) \\
		|\psi\rangle_{\text{T0}} &= \alpha|0\rangle + \beta|1\rangle \leftrightarrow \alpha\deltaE_0 + \beta\deltaE_1
	\end{align}
	\subsubsection{T0 Quantum Gates}
	Quantum gates are realized through targeted time field manipulation:
	\textbf{T0 Hadamard Gate:}
	\begin{equation}
		H_{\text{T0}} = \frac{1}{\sqrt{2}}\begin{pmatrix} 1 & 1 \\ 1 & -1 \end{pmatrix} \cdot \left(1 + \xipar \frac{\langle\deltaE\rangle}{\EPlanck}\right)
	\end{equation}
	\textbf{T0 CNOT Gate:}
	\begin{equation}
		\text{CNOT}_{\text{T0}} = \begin{pmatrix} 1 & 0 & 0 & 0 \\ 0 & 1 & 0 & 0 \\ 0 & 0 & 0 & 1 \\ 0 & 0 & 1 & 0 \end{pmatrix} \cdot \left(\mathbb{I} + \xipar \frac{\delta\Efield}{\EPlanck} \sigma_z \otimes \sigma_x\right)
	\end{equation}
	\subsection{Quantum Algorithms with T0 Improvements}
	\subsubsection{T0 Shor Algorithm}
	The factorization algorithm would be modified by the deterministic T0 evolution \textbf{[S]}:
	\begin{equation}
		P_{\text{Erfolg}}^{\text{T0}} = P_{\text{Erfolg}}^{\text{Standard}} \cdot \left(1 + \xipar \sqrt{n}\right)
	\end{equation}
	where $n$ is the bit length of the number to be factored. For RSA-2048, this estimate would give a success probability improved by $\sim 10^{-2}$.

	\textit{(Corrected on 30~Sept.~2026: previously \enquote{the number to be factored} -- only with the bit length $n = 2048$ does $\xi\sqrt{n} = 6 \times 10^{-3} \sim 10^{-2}$ hold; with the number itself ($\sim 2^{2048}$), $\xi\sqrt{n}$ would be of order $10^{304}$.)}
	\subsubsection{T0 Grover Algorithm}
	The database search would be modified by energy field focusing \textbf{[S]}:
	\begin{equation}
		N_{\text{Iterationen}}^{\text{T0}} = \frac{\pi}{4}\sqrt{N} \left(1 - \xipar \ln N\right)
	\end{equation}
	This would give a small logarithmic reduction of the iteration count for large databases.
	\section{Bell Inequalities and T0 Locality}
	\subsection{T0-Modified Bell Inequalities}
	The Bell inequalities receive small corrections through the T0 time field:
	\begin{tcolorbox}[colback=red!5!white,colframe=red!75!black,title={T0 Bell Corrections}]
		\textbf{Modified CHSH Inequality:}
		\begin{equation}
			|E(a,b) - E(a,b') + E(a',b) + E(a',b')| \leq 2 + \xipar \Delta_{\text{T0}}
		\end{equation}
		where $\Delta_{\text{T0}}$ is the time field correction:
		\begin{equation}
			\Delta_{\text{T0}} = \frac{\langle|\deltaE_A - \deltaE_B|\rangle}{\EPlanck}
		\end{equation}
	\end{tcolorbox}
	\subsection{Local Reality with T0 Fields}
	The T0 theory proposes a local-realistic interpretation of quantum correlations \textbf{[S]}:
	\subsubsection{Hidden Variable: The Time Field}
	The T0 time field acts as a local hidden variable:
	\begin{equation}
		P(A,B|a,b,\lambda_{\text{T0}}) = P_A(A|a,T_{\text{field},A}) \cdot P_B(B|b,T_{\text{field},B})
	\end{equation}
	where $\lambda_{\text{T0}} = \{T_{\text{field},A}(t), T_{\text{field},B}(t)\}$ are the local time field configurations.
	\subsubsection{Superdeterminism through T0 Correlations}
	The T0 time field leads to a superdeterministic interpretation without ''spooky action at a distance'':
	\begin{align}
		T_{\text{field},A}(t) &= T_{\text{field},\text{common}}(t-r/c) + \delta T_{\text{field},A}(t) \\
		T_{\text{field},B}(t) &= T_{\text{field},\text{common}}(t-r/c) + \delta T_{\text{field},B}(t)
	\end{align}
	The common time field history is meant to explain the correlations without violating locality. This requires that source and measurement settings are statistically dependent through this common history (superdeterminism); without that assumption, the Bell bound holds for the factorized form above.
	\section{Experimental Tests of T0 Quantum Mechanics}
	\subsection{High-Precision Interferometry}
	\subsubsection{Atom Interferometer with T0 Signatures}
	Atom interferometers could detect T0 effects through phase shifts:
	\begin{equation}
		\Delta\phi_{\text{T0}} = \frac{m \cdot v \cdot L}{\hbar} \cdot \xipar \frac{\langle\deltaE\rangle}{\EPlanck}
	\end{equation}
	For cesium atoms ($m = 2.21 \times 10^{-25}$~kg) in a 1-meter interferometer with $v = 1$--$10$~m/s:
	\begin{equation}
		\Delta\phi_{\text{T0}} \sim 2 \times 10^{-23} \ldots 2 \times 10^{-22} \text{ rad} \times \frac{\langle\deltaE\rangle}{1 \text{ eV}}
	\end{equation}
	\textit{(Corrected on 30~Sept.~2026: previously $\sim 10^{-18}$~rad without stating $v$ -- $m v L/\hbar \cdot \xi \cdot 1\,\text{eV}/E_P = 2.3 \times 10^{-23}$~rad per m/s; $10^{-18}$~rad would require $v \approx 4 \times 10^4$~m/s, far above usual atom velocities.)}
	\subsubsection{Gravitational Wave Interferometry}
	LIGO/Virgo could measure T0 corrections in gravitational wave signals:
	\begin{equation}
		h_{\text{T0}}(f) = h_{\text{GR}}(f) \left(1 + \xipar \left(\frac{f}{f_{\text{Planck}}}\right)^2\right)
	\end{equation}
	\subsection{Quantum Computer Benchmarks}
	\subsubsection{T0 Quantum Error Rate}
	For T0 quantum computers a slightly lower error rate would result \textbf{[S]}:
	\begin{equation}
		\epsilon_{\text{gate}}^{\text{T0}} = \epsilon_{\text{gate}}^{\text{Standard}} \cdot \left(1 - \xipar \frac{E_{\text{gate}}}{\EPlanck}\right)
	\end{equation}
	\section{Philosophical Implications of T0 Quantum Mechanics}
	\subsection{Determinism vs. Quantum Randomness}
	The T0 theory offers a deterministic interpretation of quantum randomness:
	\begin{tcolorbox}[colback=purple!5!white,colframe=purple!75!black,title={T0 Determinism},breakable,width=\textwidth]
		\textbf{Quantum Randomness as Effective Unpredictability:}
		What appears as fundamental randomness in standard QM is deterministic time field dynamics in the T0 theory.
		These dynamics lead to practically unpredictable, but in principle determined outcomes.
		\begin{equation}
			\begin{split}
				\text{``Randomness''} &= \text{Deterministic} \\
				&\quad \text{Time Field Evolution} \\
				&\quad + \text{Practical} \\
				&\quad \text{Unpredictability}
			\end{split}
		\end{equation}
	\end{tcolorbox}
	\subsection{Interpretation of the Measurement Problem}
	For the measurement problem of quantum mechanics, the T0 theory proposes the following route \textbf{[S]}:
	\begin{itemize}
		\item \textbf{No Collapse:} Wave functions evolve continuously
		\item \textbf{Measurement Devices:} Macroscopic T0 field configurations
		\item \textbf{Definite Outcomes:} Deterministic time field interactions
		\item \textbf{Born Rule:} Expected to emerge from T0 field dynamics (not derived here)
	\end{itemize}
	\subsection{Locality and Realism in the T0 Picture}
	The T0 theory aims to retain locality and realism; this requires the superdeterministic interpretation described above \textbf{[S]}:
	\begin{align}
		\text{Locality:} &\quad \text{All interactions mediated by local T0 fields} \\
		\text{Realism:} &\quad \text{Particles have definite properties before measurement} \\
		\text{Causality:} &\quad \text{No superluminal information transfer}
	\end{align}
	\section{Technological Applications}
	\subsection{T0 Quantum Computer Architecture}
	\subsubsection{Hardware Implementation}
	T0 quantum computers could be realized through controlled time field manipulation:
	\begin{itemize}
		\item \textbf{Time Field Modulators:} High-frequency electromagnetic fields
		\item \textbf{Energy Field Sensors:} Ultra-precise field measurement devices
		\item \textbf{Coherence Control:} Stabilization through time field feedback
		\item \textbf{Scalability:} Natural decoupling of neighboring qubits
	\end{itemize}
	\subsubsection{Quantum Error Correction with T0}
	T0-specific error correction codes:
	\begin{equation}
		|\psi_{\text{kodiert}}\rangle = \sum_i c_i |i\rangle \otimes |T_{\text{field},i}\rangle
	\end{equation}
	The time field could serve as a syndrome for error detection.
	\subsection{Precision Measurement Technology}
	\subsubsection{T0-Enhanced Atomic Clocks}
	For atomic clocks, the T0 correction is:
	\begin{equation}
		\delta f / f_0 = \delta f_{\text{Standard}} / f_0 - \xipar \frac{\Delta E_{\text{Transition}}}{\EPlanck}
	\end{equation}
	\subsubsection{Gravitational Wave Detectors}
	Possible change in sensitivity through T0 field calibration \textbf{[S]}:
	\begin{equation}
		h_{\text{min}}^{\text{T0}} = h_{\text{min}}^{\text{Standard}} \cdot \left(1 - \xipar \sqrt{f \cdot t_{\text{int}}}\right)
	\end{equation}
	\section{Standard Model Extensions}
	\subsection{T0-Extended Standard Model}
	The Standard Model is formally extended by T0 terms:
	\begin{equation}
		\mathcal{L}_{\text{SM}}^{\text{T0}} = \mathcal{L}_{\text{SM}} + \mathcal{L}_{\text{T0-Feld}} + \mathcal{L}_{\text{T0-Interaction}}
	\end{equation}
	where:
	\begin{align}
		\mathcal{L}_{\text{T0-Feld}} &= \frac{\xipar}{\EPlanck^2} (\partial \Tfield)^2 \\
		\mathcal{L}_{\text{T0-Interaction}} &= \xipar \sum_i g_i \bar{\psi}_i \gamma^\mu \partial_\mu \Tfield \psi_i
	\end{align}
	\subsection{Hierarchy Problem: T0 Approach}
	For the hierarchy problem, the T0 structure proposes the following relation \textbf{[S]}:
	\begin{equation}
		\frac{M_{\text{Planck}}}{M_{\text{EW}}} = \frac{1}{\sqrt{\xipar}} \approx \frac{1}{\sqrt{1.33 \times 10^{-4}}} \approx 87
	\end{equation}
	instead of the $10^{16}$ in the Standard Model. How this relation fits the measured ratio of the Planck and electroweak scales remains open in this early document.

	\textit{(Note of 30~Sept.~2026: $1/\sqrt{\xi} = 86.6$ is calculated correctly; the measured value is $E_P/v = 5 \times 10^{16}$, so the relation does not reproduce the scale ratio.)}
	\section{Critical Evaluation and Limitations}
	\subsection{Experimental Challenges}
	The experimental verification of the T0 theory requires:
	\begin{itemize}
		\item \textbf{Ultra-High Precision}: Measurements at the $10^{-22}$--$10^{-31}$ level
		\item \textbf{New Technologies}: T0 field-specific measurement devices
		\item \textbf{Long-Term Stability}: Consistent measurements over years
		\item \textbf{Systematic Control}: Elimination of all other effects
	\end{itemize}
	\subsection{Philosophical Implications}
	The T0 theory raises basic philosophical questions:
	\begin{itemize}
		\item \textbf{Free Will}: Is determinism compatible with human freedom of decision?
		\item \textbf{Epistemology}: How can we fully recognize the T0 reality?
		\item \textbf{Reductionism}: Are all phenomena reducible to T0 fields?
		\item \textbf{Emergence}: What role do emergent properties play?
	\end{itemize}